\section{Conclusion}\label{sec:conclusion}

We propose a method that combines both start-state and goal curricula for sparse-reward, long-horizon goal-conditioned reinforcement learning. We transfer the idea of bidirectional expansion from classic sampling-based planning~\citep{lavalle2001randomized, kuffner2000rrt} to curriculum generation, growing a set of start states outward from the goal and a set of goals outward from the initial state distribution, biasing both toward unexplored regions of the task space, and steering the two frontiers toward each other. Experimental results show that this leads to sample-efficient and robust policies on quadrupedal box climbing, real-world terrains, point-mass mazes, and manipulation tasks. Furthermore, our approach adds minimal computational overhead and thus does not require excessive training time.

Our method relies on a simulator that can be reset to arbitrary states, on rollouts of random actions, and on a bounded task space with a meaningful distance metric. For the latter, we rely on a Euclidean distance metric in this work, but future work could extend to or even learn different distance metrics, for example the observed number of time steps between two states proposed by~\citet{hartikainen2020dynamicaldistancelearningsemisupervised}. Instead of using random walk rollouts to estimate the reachability of states, future work could also explore the use of learned dynamics models or using the current policy to roll out trajectories.

